% AY2015 制御工学 第2問〔II〕（転記）
% 出典: 03_制御工学/original/03_制御工学/2015/2015se.pdf
% 要約: G1(s) と K0 を直列した G2(s)。ゲイン位相・ボード線図・ステップ応答、
%       G3(s) を帰還に持つフィードバック系での K1 決定と過渡特性の比較考察。

\section*{第2問〔II〕}

制御システムに関する次の問い (1)〜(4) に答えよ.

\begin{enumerate}[label=(\arabic*)]
  \item 伝達関数 $G_1(s)=\dfrac{1}{20s^{2}+22s+2}$ で表されるシステムと, ゲイン $K_0=2$ を
        直列結合したシステムを $G_2(s)$ とする. $G_2(s)$ のゲインならびに位相を求めよ.

  \item $G_2(s)$ のボード線図のうち, ゲイン線図を漸近線で描け. 折れ点等の特徴点を明示し,
        解答用紙内の縦軸・横軸の目盛には適宜数値を記入すること.

  \item $G_2(s)$ の単位ステップ応答を求め, 時間応答の概形を描け.

  \item $G_2(s)$ の応答特性を改善するために, 下図に示すフィードバックシステムを用いる.
        なお, $G_3(s)=s+1$ とする. このとき, フィードバックシステムの単位ステップ応答の
        定常値が, 問い (3) と同値になるようにゲイン $K_1$ を定めよ. また, システムの時定数や
        定常状態に至るまでの時間等を求めることにより, システム $G_2(s)$ と比較して過渡特性が
        どのように変化するかを考察し, 述べよ.

  \begin{center}
  \begin{tikzpicture}[>=Latex]
    \node (in)  at (0,0) {};
    \node (sum) [sum] at (1.0,0) {};
    \node (G2)  [block] at (2.8,0) {$G_2(s)$};
    \coordinate (nd) at (4.3,0);
    \node (K1)  [block] at (5.6,0) {$K_1$};
    \node (out) at (7.2,0) {};
    \node (G3)  [block] at (3.4,-1.5) {$G_3(s)$};

    \draw[->] (in)  -- (sum);
    \draw[->] (sum) -- (G2);
    \draw[->] (G2)  -- (K1);
    \draw[->] (K1)  -- (out);
    \fill (nd) circle (1.4pt);
    \draw (nd) -- (4.3,-1.5) -- (G3.east);
    \draw[->] (G3.west) -| (sum);

    \node at ($(sum.north west)+(0.02,0.10)$) {\scriptsize $+$};
    \node at ($(sum.south west)+(0.02,-0.10)$) {\scriptsize $-$};
  \end{tikzpicture}
  \end{center}
\end{enumerate}
